\documentclass[11pt]{article}
\usepackage[T1]{fontenc}
\usepackage{lmodern}
\usepackage{amsmath,amssymb,amsthm,mathtools}
\usepackage[margin=1in]{geometry}
\usepackage{hyperref}

\newtheorem{theorem}{Theorem}
\newtheorem{proposition}{Proposition}
\newtheorem{remark}{Remark}
\DeclareMathOperator{\esssup}{ess\,sup}
\DeclareMathOperator{\spec}{\sigma}
\newcommand{\R}{\mathbb{R}}
\newcommand{\Z}{\mathbb{Z}}
\newcommand{\F}{\mathcal{F}}

\title{E61: Exact Pattern-Formation Threshold in the Continuum}
\author{Mathematical proof for Manuscript 1}
\date{}

\begin{document}
\maketitle

\begin{abstract}
For the centered Swift--Hohenberg linear evolution on $\R^d$, we prove that
its spectral-instability threshold is exactly zero. We determine the
spectrum and semigroup norm, distinguish strong stability from uniform
exponential stability at criticality, characterize the unstable Fourier
annulus, and derive the corresponding finite-periodic-domain threshold.
The distinction between centered and uncentered control parameters is made
explicit. The theorem concerns linear instability, not nonlinear pattern
existence or localization.
\end{abstract}

\section{Continuum theorem}

\begin{theorem}[E61: continuum spectral threshold]\label{thm:E61}
Let $d\ge 1$, $a>0$, and $k_\star>0$. On $L^2(\R^d)$ define
\begin{equation}\label{eq:operator}
L_r = rI-a(\Delta+k_\star^2)^2,
\qquad D(L_r)=H^4(\R^d), \quad r\in\R.
\end{equation}
Then
\begin{equation}\label{eq:spectrum}
\spec(L_r)=(-\infty,r], \qquad
\|e^{tL_r}\|_{L^2\to L^2}=e^{rt}\quad (t\ge0).
\end{equation}
Consequently the critical centered growth parameter is
\begin{equation}\label{eq:threshold}
\boxed{\displaystyle
r_c=\min_{k\in\R^d}a(|k|^2-k_\star^2)^2=0.}
\end{equation}
More precisely:
\begin{enumerate}
\item If $r<0$, every $L^2$ solution decays exponentially.
\item If $r=0$, every $L^2$ solution tends strongly to zero as $t\to\infty$,
      although the semigroup is not uniformly exponentially stable.
\item If $r>0$, there exist nonzero $L^2$ data whose solutions grow
      exponentially in norm.
\end{enumerate}
The maximum Fourier growth rate is $r$, attained as a value of the
symbol precisely on $|k|=k_\star$.
\end{theorem}

\begin{proof}
We use the unitary Fourier transform, so
$\widehat{\Delta u}(k)=-|k|^2\widehat u(k)$. Hence $L_r$ is unitarily
equivalent to multiplication by the real-valued function
\begin{equation}\label{eq:symbol}
\lambda_r(k)=r-a(|k|^2-k_\star^2)^2.
\end{equation}
For every $k\in\R^d$,
\[
a(|k|^2-k_\star^2)^2\ge0,
\qquad\lambda_r(k)\le r.
\]
Since $k_\star>0$, the sphere
$S_{k_\star}=\{k\in\R^d:|k|=k_\star\}$ is nonempty, and
$\lambda_r(k)=r$ for each $k\in S_{k_\star}$. It follows that
\begin{equation}\label{eq:extremum}
\sup_{k\in\R^d}\lambda_r(k)=r,
\qquad \min_{k\in\R^d}a(|k|^2-k_\star^2)^2=0.
\end{equation}

The function $\lambda_r$ is continuous, has range $(-\infty,r]$, and
tends to $-\infty$ as $|k|\to\infty$. Its essential range is also
$(-\infty,r]$: every neighborhood of a value in this range has a
preimage of positive Lebesgue measure (including neighborhoods of the
maximum $r$). The spectrum of an unbounded self-adjoint multiplication
operator is its essential range, giving the first identity in
\eqref{eq:spectrum}. The Fourier formula for the semigroup is
\begin{equation}\label{eq:semigroup}
\widehat{e^{tL_r}u_0}(k)
=e^{t\lambda_r(k)}\widehat u_0(k), \qquad t\ge0.
\end{equation}
For $t\ge0$, its operator norm is the essential supremum of the
multiplier:
\[
\|e^{tL_r}\|_{2\to2}
=\esssup_{k\in\R^d}e^{t\lambda_r(k)}=e^{rt}.
\]
In particular, if $r<0$, then
\[
\|e^{tL_r}u_0\|_2\le e^{rt}\|u_0\|_2\longrightarrow0
\qquad(t\to\infty).
\]

At $r=0$, Plancherel's identity gives
\begin{equation}\label{eq:critical-decay}
\|e^{tL_0}u_0\|_2^2
=\int_{\R^d}
 e^{-2at(|k|^2-k_\star^2)^2}|\widehat u_0(k)|^2\,dk.
\end{equation}
The shell $S_{k_\star}$ has $d$-dimensional Lebesgue measure zero
(including when $d=1$, where it consists of two points). The multiplier
in \eqref{eq:critical-decay} therefore tends to zero for almost every
$k$, while remaining bounded by $1$. Since
$|\widehat u_0|^2\in L^1(\R^d)$, dominated convergence gives
\[
\lim_{t\to\infty}\|e^{tL_0}u_0\|_2=0.
\]
Nevertheless $\|e^{tL_0}\|_{2\to2}=1$ at every $t\ge0$,
so the semigroup is not uniformly exponentially stable. Note that
the shell vectors are not nonzero $L^2$ eigenfunctions; the value zero
is a spectral edge, not an $L^2$ point eigenvalue.

For $r>0$, choose $k_0\in S_{k_\star}$. By continuity of the symbol,
there is a bounded open ball $U$ centered at $k_0$ such that
$\lambda_r(k)>r/2$ on $U$. Choose nonzero $u_0\in L^2$ with
$\widehat u_0$ supported in $U$. Then
\[
\|e^{tL_r}u_0\|_2^2
=\int_U e^{2t\lambda_r(k)}|\widehat u_0(k)|^2\,dk
\ge e^{rt}\|u_0\|_2^2,
\]
and therefore $\|e^{tL_r}u_0\|_2\ge e^{rt/2}\|u_0\|_2$.
Thus instability first occurs as $r$ crosses zero, proving
\eqref{eq:threshold} and the stated classification.
\end{proof}

\section{Unstable spectral band and wavelength}

\begin{proposition}[Unstable annulus]\label{prop:band}
For $r>0$, the unstable wavevectors are exactly
\begin{equation}\label{eq:unstable}
\bigl\{k\in\R^d:\bigl||k|^2-k_\star^2\bigr|<\sqrt{r/a}\bigr\}.
\end{equation}
If $0<r<ak_\star^4$, this is the annulus
\begin{equation}\label{eq:annulus}
\sqrt{k_\star^2-\sqrt{r/a}}<|k|
<\sqrt{k_\star^2+\sqrt{r/a}}.
\end{equation}
The fastest growth occurs at $|k|=k_\star$, corresponding to
$\lambda_\star=2\pi/k_\star$. As $r\downarrow0$, the annulus collapses
onto the shell $|k|=k_\star$.
\end{proposition}

\begin{proof}
The inequality $\lambda_r(k)>0$ is equivalent to
$a(|k|^2-k_\star^2)^2<r$, which is precisely
\eqref{eq:unstable}. When $0<r<ak_\star^4$,
both bounds on $|k|^2$ are positive, so taking their square roots gives
\eqref{eq:annulus}. The maximum is characterized in
\eqref{eq:extremum}. A plane wave of wavenumber magnitude $k_\star$
has wavelength $2\pi/k_\star$. Both radii in \eqref{eq:annulus}
converge to $k_\star$ as $r\downarrow0$.
\end{proof}

\begin{remark}[Near-onset width]
Writing $|k|=k_\star+\delta$ gives
\[
\lambda_r(k)
=r-a(2k_\star\delta+\delta^2)^2
=r-4ak_\star^2\delta^2+O(\delta^3).
\]
Thus the unstable half-width is asymptotic to
$(2k_\star)^{-1}\sqrt{r/a}$ as $r\downarrow0$.
\end{remark}

\section{Finite periodic domains}

\begin{proposition}[Periodic-box threshold]\label{prop:torus}
On the flat torus $\mathbb{T}_L^d=(\R/L\Z)^d$, the linear growth rates
are
\[
\lambda_n(r)=r-a\left(\left|\frac{2\pi n}{L}\right|^2-k_\star^2\right)^2,
\qquad n\in\Z^d.
\]
The threshold is
\begin{equation}\label{eq:torus-threshold}
r_c(L)=a\min_{n\in\Z^d}
\left(\left|\frac{2\pi n}{L}\right|^2-k_\star^2\right)^2\ge0.
\end{equation}
Equality holds if and only if an admissible Fourier wavevector lies on
$|k|=k_\star$. Moreover,
\begin{equation}\label{eq:limit}
\lim_{L\to\infty}r_c(L)=0.
\end{equation}
\end{proposition}

\begin{proof}
The periodic Fourier basis has wavevectors $k_n=2\pi n/L$,
so each Fourier coefficient solves
$\partial_t\widehat u_n=\lambda_n(r)\widehat u_n$.
The first mode reaches nonnegative growth when $r$ reaches the minimum
in \eqref{eq:torus-threshold}. The minimum exists because the
nonnegative expression tends to infinity with $|n|$. It vanishes
exactly when $|2\pi n/L|=k_\star$ for some $n\in\Z^d$.

For the limit, choose integers $m(L)$ satisfying
$|m(L)-Lk_\star/(2\pi)|\le1/2$ and let
$n(L)=(m(L),0,\ldots,0)$. Then
\[
\frac{2\pi n(L)}{L}\longrightarrow k_\star e_1.
\]
By the minimum property,
\[
0\le r_c(L)
\le a\left(\left|\frac{2\pi n(L)}L\right|^2-k_\star^2\right)^2
\longrightarrow0.
\]
The squeeze theorem proves \eqref{eq:limit}.
\end{proof}

\section{Centered versus uncentered parameters}

\begin{proposition}[Control-parameter translation]\label{prop:shift}
Suppose the original (uncentered) dispersion relation is
\[
\lambda(k)=r_0+\alpha|k|^2-b|k|^4,
\qquad \alpha,b>0.
\]
Writing $k_\star^2=\alpha/(2b)$ and
$\mu=r_0+\alpha^2/(4b)$ yields
\begin{equation}\label{eq:shift}
\lambda(k)=\mu-b(|k|^2-k_\star^2)^2.
\end{equation}
Therefore the threshold is $\mu_c=0$ in centered coordinates, and
\begin{equation}\label{eq:uncentered}
\boxed{r_{0,c}=-\frac{\alpha^2}{4b}}
\end{equation}
in the original coordinates.
\end{proposition}

\begin{proof}
Expanding the square in \eqref{eq:shift},
\[
\mu-b\left(|k|^4-2k_\star^2|k|^2+k_\star^4\right)
=\mu-b|k|^4+\alpha|k|^2-\frac{\alpha^2}{4b}
=r_0+\alpha|k|^2-b|k|^4.
\]
Theorem~\ref{thm:E61} gives $\mu_c=0$, which is equivalent to
\eqref{eq:uncentered}.
\end{proof}

\section{Scope of the result}

Theorem~\ref{thm:E61} proves a \emph{linear spectral-instability
threshold}. It does not assert nonlinear saturation, existence or
uniqueness of patterns, orbital stability of nonzero equilibria,
or localized finite-energy bound states. In particular, although an
unprojected nonlinear Swift--Hohenberg evolution such as
\[
\partial_t A=rA-a(\Delta+k_\star^2)^2A-\beta|A|^2A
\]
has the same linearization at $A=0$, the existence and stability of
nonzero patterns require additional arguments. Nor does the present
proof derive the linear operator from any more fundamental WCT model.

\end{document}